%% 7_conclusion
\section{Conclusion}
\label{sec:conclusion}

Architectures for radio field reconstruction have been compared against each other
for several years without the comparison revealing what actually limits them. Under
one protocol, nine learned methods spanning convolutional, cascaded, state-space,
transformer, adversarial and diffusion designs vary over $20.1\%$ in line-of-sight
error and only $9.6\%$ in non-line-of-sight error, which is the regime that
dominates the total. A plateau this indifferent to design is a statement about the
inputs, not the models.

The missing input is a line integral. Supplying the fraction of the
transmitter-to-pixel ray that lies inside building material costs $432$ parameters
in a single network, and under $600$ in the cascade, yet moves the result past
every method in the comparison, including in a plain
single network that has none of the architectural machinery its competitors carry.
Geometry and depth turn out to be substitutes rather than complements, and a fifth
of their combined expected benefit is lost to overlap, which is a reason to spend
effort on conditioning rather than on capacity.

The same finding explains why this literature transfers poorly. The line integral
is anchored at the transmitter, and indoor cellular crowdsensing locates the
transmitter in fewer than one cell in five. Marginalising the anchor away destroys
the field it was supposed to preserve, and on real measurements classical
interpolation remains ahead of every learned reconstructor we trained. The gap
between simulated and deployed performance is therefore not a domain shift to be
closed by fine-tuning. It is an input that is present in one setting and absent in
the other.

Direction is the quantity worth recovering next. Inferring where energy arrives
from, using the measurements themselves rather than assuming a known source or
discarding direction altogether, is the step that would make geometry conditioning
deployable. The failure in \cref{sec:txfree_result} is specific and diagnosable,
which makes it a starting point rather than a wall.
